
\chapter{כלכלת אסימונים (\textenglish{Token Economics})}

\section{מודל עלות הסוכן}

ד"ר סגל \cite{segal2024orchestration} הציע נוסחה מצטברת לעלות אסימון בסוכן:

\begin{equation}
WC_n = WC_{n-1} + Q_n + R_n + A_n
\end{equation}

כאשר \(WC_n\) הוא סך האסימונים בסבב \textenglish{n}, \(Q_n\) אסימוני השאלה, \(R_n\) אסימוני התוצאה ו-\(A_n\) אסימוני פעולת הסוכן.

\section{ניתוב חכם לחיסכון}

ניתוב כישורים דינמי (\textenglish{Router-Skill}) מפחית \(WC_n\) על ידי טעינת כלים רק בעת הצורך:

\begin{equation}
\Delta WC = \sum_{n=1}^{N} \mathbf{1}[\text{skill not needed}_n] \cdot C_{\text{skill load}}
\end{equation}

ניסויים הראו חיסכון של \textenglish{23\%} בעלות הכוללת עבור צינורות עם \textenglish{$> 10$} כישורים ושימוש ממוצע \textenglish{$< 40\%$} לכישור \cite{segal2024orchestration}.

\section{מקבוליות ועלות}

\begin{table}[H]
\centering
\caption{השוואת מסגרות תיאום: \textenglish{Latency} מול \textenglish{Token Cost}}
\label{tab:frameworks}
\begin{english}
\begin{tabular}{lccc}
\toprule
\textbf{Framework} & \textbf{Latency (s)} & \textbf{Token Cost} & \textbf{Error Rate} \\
\midrule
ReAct (single)   & 12.4 & 1.0× (baseline) & 18\% \\
AutoGen          & 18.7 & 2.3×             & 11\% \\
MetaGPT          & 31.2 & 4.1×             & 7\%  \\
CrewAI + Router  & \textbf{14.1} & \textbf{1.4×} & \textbf{6\%} \\
\bottomrule
\end{tabular}
\end{english}
\end{table}
